% ============================================================
% SECTION 18 — REQUEST FLOWS
% ============================================================

\section{Request Flows}

\begin{tcolorbox}[summarybox={\iconFlow[1.5] Step-by-Step Execution Lifecycle}]
\color{textdark}
To understand how the decoupled layers (FastAPI, Redis, Celery, PostgreSQL) interact in practice, it is necessary to trace the exact sequence of events for the system's core operations.

This section maps the two most critical paths: initiating a payment recovery and initiating a refund.
\end{tcolorbox}

\vspace{0.8cm}

% --- THE PAYMENT RECOVERY FLOW ---
\subsection{Flow 1: Initiating a Payment Recovery}

This is the primary business function of RecoverAI V2: attempting to salvage a failed transaction. Note how the HTTP request terminates long before the actual financial execution occurs.

\begin{center}
\begin{tikzpicture}[
    node distance=0.8cm and 1.8cm,
    actor/.style={rectangle, draw=primary, thick, fill=bglight, minimum width=2.5cm, minimum height=0.8cm, rounded corners=4pt, font=\small\bfseries, align=center},
    db/.style={cylinder, draw=accent, thick, fill=accent!10, shape border rotate=90, aspect=0.25, minimum width=2cm, minimum height=1.2cm, font=\small\bfseries, align=center},
    msg/.style={rectangle, draw=textdark, dashed, minimum width=3cm, minimum height=0.6cm, font=\scriptsize, align=center},
    arrow/.style={-{Stealth[scale=1.1]}, thick, draw=textdark}
]

    % Actors
    \node[actor] (client) {Client};
    \node[actor, below=1.2cm of client] (api) {FastAPI};
    \node[actor, below=1.2cm of api] (redis) {Redis Broker};
    \node[actor, below=1.2cm of redis] (worker) {Celery Worker};
    \node[actor, below=1.2cm of worker] (gw) {Gateway};
    \node[db, right=3cm of api, yshift=-1.5cm] (db) {PostgreSQL};

    % 1. Client to API
    \draw[arrow] (client.south) -- node[left, font=\scriptsize] {1. POST \code{/recover} + API Key + Idemp Key} (api.north);
    
    % 2. API to DB (Idempotency)
    \draw[arrow] (api.east) -- node[above, font=\scriptsize, sloped] {2. Insert IdempotencyRecord} (db.120);
    
    % 3. API to DB (Transaction)
    \draw[arrow] (api.south east) -- node[above, font=\scriptsize, sloped] {3. Insert Transaction (PENDING)} (db.150);
    
    % 4. API to Redis
    \draw[arrow] (api.south) -- node[left, font=\scriptsize] {4. Enqueue \code{process\_orchestrator}} (redis.north);
    
    % 5. API to Client (Return)
    \draw[arrow, draw=success] (api.west) to[out=180, in=180] node[left, font=\scriptsize, text=success] {5. 202 Accepted} (client.west);
    
    % 6. Redis to Worker
    \draw[arrow] (redis.south) -- node[left, font=\scriptsize] {6. Consume Task} (worker.north);
    
    % 7. Worker to DB
    \draw[arrow] (worker.east) -- node[below, font=\scriptsize, sloped] {7. OCC Update to AUTHORIZED} (db.210);
    
    % 8. Worker to Gateway
    \draw[arrow] (worker.south) -- node[left, font=\scriptsize] {8. \code{ExecutionGuard.execute()}} (gw.north);
    
    % 9. Gateway to Worker
    \draw[arrow, draw=success] (gw.east) to[out=0, in=0] node[right, font=\scriptsize, text=success] {9. Success (200 OK)} (worker.east);
    
    % 10. Worker to DB (Final)
    \draw[arrow] (worker.north east) -- node[below, font=\scriptsize, sloped] {10. Final Update (SUCCEEDED)} (db.240);

\end{tikzpicture}
\end{center}

\vspace{0.8cm}

% --- THE REFUND FLOW ---
\subsection{Flow 2: Initiating a Refund}

Unlike recoveries (which require LLM diagnosis and policy evaluation in the background), refunds are synchronous. The API holds the connection open, directly calls the \code{RefundService}, and waits for the gateway to respond.

\begin{center}
\begin{tikzpicture}[
    node distance=0.8cm and 1.8cm,
    actor/.style={rectangle, draw=secondary, thick, fill=bglight, minimum width=2.5cm, minimum height=0.8cm, rounded corners=4pt, font=\small\bfseries, align=center},
    db/.style={cylinder, draw=accent, thick, fill=accent!10, shape border rotate=90, aspect=0.25, minimum width=2cm, minimum height=1.2cm, font=\small\bfseries, align=center},
    arrow/.style={-{Stealth[scale=1.1]}, thick, draw=textdark}
]

    % Actors
    \node[actor] (client) {Client};
    \node[actor, below=1.2cm of client] (api) {FastAPI Router};
    \node[actor, below=1.2cm of api] (service) {Refund Service};
    \node[actor, below=1.2cm of service] (gw) {Gateway};
    \node[db, right=3cm of api, yshift=-1cm] (db) {PostgreSQL};

    % 1. Client to API
    \draw[arrow] (client.south) -- node[left, font=\scriptsize] {1. POST \code{/refund} + Idemp Key} (api.north);
    
    % 2. API to DB (Idempotency)
    \draw[arrow] (api.east) -- node[above, font=\scriptsize, sloped] {2. Insert IdempotencyRecord} (db.120);
    
    % 3. API to Service
    \draw[arrow] (api.south) -- node[left, font=\scriptsize] {3. \code{initiate\_refund(txn\_id)}} (service.north);
    
    % 4. Service to DB (OCC)
    \draw[arrow] (service.east) -- node[below, font=\scriptsize, sloped] {4. \code{UPDATE ... WHERE status='NONE'}} (db.210);
    
    % 5. Service to Gateway
    \draw[arrow] (service.south) -- node[left, font=\scriptsize] {5. Gateway API Call} (gw.north);
    
    % 6. Gateway to Service
    \draw[arrow, draw=success] (gw.east) to[out=0, in=0] node[right, font=\scriptsize, text=success] {6. Gateway returns Processing} (service.east);
    
    % 7. Service to DB (Update)
    \draw[arrow] (service.north east) -- node[below, font=\scriptsize, sloped] {7. Update to \code{REFUND\_PROCESSING}} (db.240);
    
    % 8. Service to API
    \draw[arrow] (service.west) to[out=180, in=180] node[left, font=\scriptsize] {8. Return Model} (api.west);
    
    % 9. API to Client
    \draw[arrow, draw=success] (api.north west) to[out=135, in=225] node[left, font=\scriptsize, text=success] {9. 200 OK} (client.south west);

\end{tikzpicture}
\end{center}

\begin{tcolorbox}[infobox={Synchronous vs. Asynchronous Refunds}]
Notice that the Refund Flow ends in \code{REFUND\_PROCESSING}, not \code{REFUNDED}. Because actual fund transfers take days to settle between banks, the synchronous API call only guarantees that the refund was successfully \textit{submitted} to the gateway. The system relies entirely on the Webhook Architecture (Section 13) to receive the final \code{refund.processed} or \code{refund.failed} event days later.
\end{tcolorbox}
